\documentclass[12pt, a4paper]{article}
\usepackage[T1]{fontenc}
\usepackage{lmodern}
\usepackage{amsmath, amsthm, amssymb, mathtools}
\usepackage[margin=1in]{geometry}
\usepackage{xr-hyper}
\usepackage[colorlinks=true,linkcolor=blue,citecolor=blue]{hyperref}
\usepackage{cleveref}
\usepackage{graphicx, subcaption}
\usepackage{natbib}
\usepackage{booktabs}
\usepackage{enumerate}

% Code references: scaled monospace with line-breaking
\newcommand{\code}[1]{\texttt{\small #1}}

\newtheorem{theorem}{Theorem}[section]
\newtheorem{lemma}[theorem]{Lemma}
\newtheorem{proposition}[theorem]{Proposition}
\newtheorem{corollary}[theorem]{Corollary}
\newtheorem{definition}{Definition}[section]
\newtheorem{remark}{Remark}[section]
\newtheorem{conjecture}[theorem]{Conjecture}
\newtheorem{derivation}[theorem]{Derivation}
\newtheorem{observation}[theorem]{Observation}
\numberwithin{equation}{section}

\newcommand{\dd}{\mathrm{d}}
\newcommand{\seriestitle}{Why Rotating Fluids Make Polygons}


\externaldocument[I-]{../paper-1-mathematics/main}
\externaldocument[II-]{../paper-2-physics/main}
\externaldocument[III-]{../paper-3-gravity/main}
\externaldocument[IV-]{../paper-4-field-theory/main}
\externaldocument[V-]{../paper-5-cosmology/main}

\title{\seriestitle\\[6pt]\large Paper VI: Discussion, Predictions, and Status}
\author{Gordon Speagle}
\date{}

\begin{document}
\maketitle

\begin{abstract}
This paper provides the series-level audit: parameter accounting,
status classification (Theorem / Derivation / Conjecture),
cumulative uncertainty analysis, and falsifiable predictions
for the Havelock field theory developed in Papers~I--V.
From one discrete input ($N = 7$), one dimensionful scale ($M_P$),
and three structural identifications
(Seifert geometry, BF threshold, minimal CS levels),
the series obtains the Standard Model gauge group, the Weinberg
angle, the CKM phase, three generations, the cosmological
constant, the Hubble parameter, the dark matter fraction,
and the baryon asymmetry.
Six binary predictions are testable within two decades.
\end{abstract}

\section{Series overview}
\label{sec:overview}

The Havelock field theory obtains the Standard Model and
general relativity from the stability of point-vortex polygons.
The derivation chain is:
\begin{enumerate}
\item \emph{Paper~I (Mathematics):}
  The Havelock decomposition
  $\lambda_m = C_1(\text{surface}) - m(N{-}m)/2$;
  the sign rule; algebraic-integer thresholds;
  the constrained energy minimum for $N \le 7$.
\item \emph{Paper~II (Physics):}
  The quartic normal form at $N = 7$;
  the blob and Rossby bridges;
  BEC and planetary predictions.
\item \emph{Paper~III (Gravity):}
  The CMS--CS Casimir identification;
  the graviton at $N = 7$;
  the vacuum Einstein equations from polygon entropy
  maximisation; the Onsager--WDW correspondence.
\item \emph{Paper~IV (Gauge theory):}
  The Seifert manifold
  $\mathbf{H}^2 \times_N S^1$;
  $\mathrm{SU}(3) \times \mathrm{SU}(2) \times \mathrm{U}(1)$
  from the McKay correspondence and CS gravity;
  $\sin^2\theta_W = 3/11$; the CKM phase $\delta = 70.2^\circ$;
  the fermion mass hierarchy.
\item \emph{Paper~V (Cosmology):}
  The $N = 11$ instanton;
  $H_0 = 67.4$\,km/s/Mpc; $\Omega_\Lambda = 68.5\%$;
  dark matter from frozen KK modes; baryogenesis;
  neutrino masses; radion inflation.
\end{enumerate}


\section{Parameter accounting}
\label{sec:parameters}

\paragraph{Summary.}
One discrete input ($N = 7$), one dimensionful scale ($M_P$),
and three structural identifications---each argued to be the
unique consistent option---produce the seventeen testable outputs
enumerated in \S\ref{sec:status}.
No continuous free parameters are adjusted.

\paragraph{Structural framework.}
Each structural element is the \emph{unique} option
compatible with the requirements of the vortex system:
\begin{enumerate}
\item \emph{Seifert manifold.}
  Uniquely determined by the Thurston classification
  (Proposition~\ref*{IV-prop:spacetime}).
\item \emph{BF threshold $\mu_4 = 1$.}
  The $\mathbf{H}^2$ Breitenlohner--Freedman bound gives
  the scalar spectrum $M^2 = \Delta(\Delta{-}1) \ge -1/4$.
  The first KK mode above the threshold
  ($\Delta = 3/2$, $M^2 > 0$) has $\mu_4 = \Delta - 1/2 = 1$.
  This is determined by the $\mathbf{H}^2$ spectrum.
\item \emph{CS levels $k = K = 1$.}
  For compact groups: $k \in \mathbb{Z}^+$
  (from $\pi_3(G) = \mathbb{Z}$).
  At $k = 1$: the integrable representations of SU$(2)_1$
  are $j = 0$ and $j = 1/2$ --- the Higgs doublet is the
  maximal representation (at $k \ge 2$ there would be
  additional states not present in the SM).
  At $K = 1$: the compact boson is at the self-dual radius
  (the unique T-duality fixed point with enhanced SU$(2)$
  symmetry).
\end{enumerate}
Within this structure:

\paragraph{Input.}
The single discrete input is $N = 7$ (the polygon number),
determined by the Pell equation (Theorem~\ref*{V-thm:n-selection}(b)).
The cosmological polygon number $N = 11$ follows from Chern class
additivity (Theorem~\ref*{V-thm:n-selection}(c)).

\paragraph{Derived from~$N$ (no freedom).}
The central charge $c = 12\,b(N)$, CS phase~$\theta_{\mathrm{CS}}$,
Pell threshold $\xi^* = 8 {-} 3\sqrt{7}$,
Weinberg angle $\sin^2\theta_W = 3/11$,
$3$~generations $= (N{-}1)/2$,
$4$~texture zeros (from $\mathbb{Z}_7$ charge conservation),
Wolfenstein hierarchy,
and CKM phase $\delta = \tfrac{1}{2}\log\cosh(\pi) = 70.2^\circ$.

\paragraph{Derived from the BF threshold and $\mathrm{SU}(2)$ representation theory.}
The fermion bulk mass on the $S^1$ fiber is
$\mu_4 = \mu_{4,\mathrm{bare}} \mp T_3$, where:
\begin{itemize}
\item $\mu_{4,\mathrm{bare}} = 1$: in the KK reduction, the fiber mass
  $\mu_4$ maps to the bulk conformal dimension $\Delta$ via $\Delta = \mu_4 + 1/2$
  (the standard mass--dimension relation for Dirac fields on
  AdS$_3$).
  The BF threshold $M^2 = \Delta(\Delta{-}1) = 0$ at $\Delta = 1$ then gives
  $\mu_4 = \Delta - 1/2 = 1/2$ for the \emph{massless} mode, and
  $\mu_{4,\mathrm{bare}} = 1$ for the first massive level
  (the unique stability boundary on~$\mathbf{H}^2$);
\item $T_3 = \pm\tfrac{1}{2}$ (the weak isospin of the Higgs doublet
  component in the Yukawa coupling).
\end{itemize}
This gives $\mu_4^{\mathrm{up}} = 1 - \tfrac{1}{2} = \tfrac{1}{2}$
(conformal profile, heaviest generation) and
$\mu_4^{\mathrm{dn}} = 1 + \tfrac{1}{2} = \tfrac{3}{2}$.
The $k = K = 1$ reasoning is as above.

\paragraph{Neutrino sector.}
The Galois centrifugal suppression $a = 1/N = 1/7$ is derived from the
orbifold angular scale.
The Dirac KK mass scale $M_{\mathrm{KK}}$ is set by eq.~\eqref{V-eq:M-R}
($M_{\mathrm{KK}} = M_{\mathrm{poly}} \times \varepsilon_7^7 \times 6$);
the overall normalisation is calibrated to
$\Delta m^2_{\mathrm{atm}}$, making this one input.
The \emph{ratio} $\Delta m^2_{\mathrm{atm}}/\Delta m^2_{\mathrm{sol}}$
and the mass hierarchy are genuine predictions
($M_{\mathrm{KK}}$-independent).

\paragraph{Total.}
One discrete input ($N = 7$) plus one dimensionful
scale ($v$ or $M_P$, needed to set units in any theory),
plus the structural identifications (Thurston $\widetilde{\mathrm{SL}(2,\mathbb{R})}$
geometry, BF-threshold identification with the $T_3$ split,
$k = K = 1$ for the compact gauge levels),
yielding $17$~testable outputs ($15$~predictions and
$2$~consistency checks in the master table), of which
$11$ are independent after accounting for Yukawa-texture
correlations and CKM unitarity.
The Standard Model, by comparison, has $19$~free parameters
for roughly the same set of observables.
The structural choices above are determined by the geometry
and representation theory, but differ from quantities
computed by integration (WKB actions, Legendre functions).

\paragraph{Sensitivity to structural modifications.}
Each structural choice, if altered, breaks an
\emph{internal} prediction (computed from the
framework's own formulae, then compared to data):
\begin{itemize}
\item Seifert $\to$ product ($e = 0$):
  no parity violation, $\mathrm{SU}(2)_L = \mathrm{SU}(2)_R$
  (contradicts maximal parity violation in weak decays).
\item BF threshold $\mu_4 \ne 1$:
  at $\mu_4 = 0$ (massless): all generations degenerate
  (contradicts the fermion mass hierarchy);
  at $\mu_4 = 2$: $c_3 = 2$ gives $m_b/m_t \sim 10^{-4}$
  (too small by $10^2$).
\item CS levels $k \ge 2$:
  the integrable representations of $\mathrm{SU}(2)_k$
  include $j = 1, 3/2, \ldots$ (extra states absent
  from the SM spectrum).
\item $N \ne 11$:
  $\Lambda_3 = (N^2{-}16)/16$ gives the wrong $H_0$
  for any $N \ne 11$
  ($N_{\mathrm{total}} = 10$: $N_{\mathrm{grav}} = 6$ is stable,
  no instanton, $H_0 = 0$;
  $N_{\mathrm{total}} = 12$: $N_{\mathrm{grav}} = 8$ gives
  $S_{\mathrm{BO}}(8) = 31.8 \gg S_{\mathrm{BO}}(7) = 18.3$,
  collapsing the hierarchy;
  even keeping the hierarchy fixed,
  $\Lambda_3(12) = 8$ gives $H_0 = 74$\,km/s/Mpc,
  excluded at ${>}\,5\sigma$).
\end{itemize}
The structural identifications have no continuous freedom:
each choice is the unique option consistent with
observation.


\section{Discussion}
\label{sec:discussion}

This section classifies every result by logical status
(\S\ref{sec:status}), propagates uncertainties through
the derivation chain (\S\ref{sec:uncertainty}), compares
to other zero-parameter approaches (\S\ref{sec:comparison}),
and identifies falsifiable predictions (\S\ref{sec:falsifiable}).

\subsection{Status classification}
\label{sec:status}
Three levels of logical status are distinguished:
\begin{center}
\small
\begin{tabular}{lp{8cm}}
\toprule
Status & Meaning \\
\midrule
\textbf{Theorem} & Proved from the theory's hypotheses
  (including the vortex-gravity identification) using
  established mathematics (Havelock identity, McKay
  correspondence, Witten CS/WZW, Feynman--Kac). \\
\textbf{Derivation} & Follows from the theory via a
  chain of identified steps, each with controlled
  approximations (WKB, Markovian, Born--Oppenheimer).
  The result depends on the structural identifications
  (Seifert geometry, $N = 7$). \\
\textbf{Conjecture} & Numerically verified but
  the derivation chain has at least one step that is
  physically motivated rather than mathematically proved. \\
\bottomrule
\end{tabular}
\end{center}

\begin{center}
\small
\begin{tabular}{lllc}
\toprule
Result & Role & Status & Match \\
\midrule
\multicolumn{4}{l}{\emph{Structural (forced by $N = 7$ and $M_P$)}} \\
$\mathrm{SU}(3)$ representation & Derived & Theorem & --- \\
$\mathrm{SU}(3)$ gauge dynamics & Derived & Derivation$^\dagger$ & --- \\
$\mathrm{SU}(2) \times \mathrm{U}(1)$ & Derived & Derivation & --- \\
$3$ generations $= (N{-}1)/2$ & Derived & Derivation & --- \\
$4$D graviton (spin $2$) & Derived & Derivation & --- \\
Vacuum Einstein + full Riemann & Derived & Derivation & --- \\
WDW equation from stat.\ mech. & Derived & Derivation & --- \\
$M_P/M_{\mathrm{EW}}$ hierarchy & Derived & Derivation & $0.003\%$ \\
\midrule
\multicolumn{4}{l}{\emph{Predictions (computed before comparison)}} \\
$\delta_{\mathrm{CKM}} = 70.2^\circ$ & Prediction & Derivation & $0.4\sigma$ \\
$\sin^2\theta_W = 3/11$ & Prediction & Derivation$^*$ & $3.9\%$ \\
$|V_{cb}| = 0.044$ & Prediction & Derivation & $6\%$ \\
$m_s/m_b = 0.023$ & Prediction & Derivation$^\dagger$ & $4\%$ \\
$H_0 = 67.4$ km/s/Mpc & Prediction & Derivation & $2 \times 10^{-5}$ \\
$\Omega_{\mathrm{DM}} = 26.9\%$ & Prediction & Conjecture & $0.4\sigma$ \\
$\Omega_b = 5.0\%$ & Prediction & Derivation & $0.6\sigma$ \\
$\eta_B = 2 \times 10^{-9}$ & Prediction & Derivation & $\times 0.8$--$3$ \\
$\Delta m^2$ ratio $= 31.5$ & Prediction & Derivation & $3\%$ \\
$m_H = 125 \pm 3$ GeV & Prediction & Derivation & $2\%$ \\
$\theta_{\mathrm{QCD}} = 0$ & Prediction & Derivation & --- \\
Normal hierarchy, $\Sigma m_\nu = 0.059$\,eV & Prediction & Derivation & JUNO \\
$w = -1 + O(1/c_{11}^2)$ & Prediction & Derivation & DESI \\
No $0\nu\beta\beta$ (Dirac) & Prediction & Derivation & LEGEND \\
Proton exactly stable & Prediction & Derivation & Hyper-K \\
EWPT GW signal, $\Omega h^2 \sim 10^{-9}$ & Prediction & Derivation & LISA \\
\midrule
\multicolumn{4}{l}{\emph{Confinement}} \\
$\alpha_s(M_Z) = 0.115$--$0.122$ & Prediction & Derivation & $0.1180 \pm 0.0009$ \\
$\sigma/\Lambda^2 = 5.97$ & Prediction & Derivation & $6.25 \pm 0.5$ \\
$\langle P\rangle = 0$ (confined) & Structural & Theorem & --- \\
\midrule
\multicolumn{4}{l}{\emph{Consistency checks}} \\
Coupling lock $\alpha/(8\pi G) = 1/(2\pi^2)$ & Check$^\ddagger$ & Identity & exact \\
Einstein + matter (Lashkari) & Check & Derivation & --- \\
\bottomrule
\end{tabular}

\smallskip
\footnotesize
Match column: \% = deviation from observation;
$\sigma$ = standard deviations;
$\times$ = multiplicative factor;
named experiment = future test;
--- = no empirical target.
$\dagger$ = the DHVW orbifold at irrational $c = 51.57$
is proved consistent (\S\ref*{IV-sec:su3-mckay}: unitarity,
modular invariance, OPE closure, positivity).
Non-integrable primaries contribute $1.2\%$ to the gauge
current four-point function (computed from the Zamolodchikov
conformal block decomposition; see Code Availability),
below the geometric-tier accuracy of all predictions.
\normalsize
\end{center}
\noindent
${}^*$The Weinberg angle uses the Lie-algebra Casimir
$C_2(j{=}1) = 2$ at the critical Havelock mode of $N = 4$.
The integrable bound $j \le 1/2$ of
$\widehat{\mathfrak{su}}(2)_1$ constrains TQFT states,
not perturbative gauge couplings
(\S\ref*{IV-sec:weinberg}).\\[2pt]
${}^\ddagger$The coupling lock is an algebraic identity
within the formalism ($\alpha$ and $G$ both derive from~$c$);
it is not independently testable at the CS threshold.\\[2pt]
${}^\dagger$The mass ratio $m_s/m_b$ is computed from the
BF instanton (\S\ref*{IV-sec:mass-hierarchy}):
$m_s/m_b = (v/M_{\mathrm{poly}})^{2\Delta c}/K^2$,
where $\Delta c = 0.351$ (including the $O(1/c)$ Euler
class and BO corrections;
eq.~\eqref*{IV-eq:Dc-corrected}) and $K = 0.548$.
At two-loop ($M_{\mathrm{poly}} = 294$\,TeV): $0.023$
(observed $0.024$, $4\%$ match).

\subsection{Cumulative uncertainty}
\label{sec:uncertainty}
The derivation chain from polygon stability to the Standard
Model passes through three distinct types of link:
\begin{enumerate}
\item \emph{Theorem links} (Havelock identity, McKay
  correspondence, CS/WZW, Liouville--Arnold):
  mathematically proved.  No propagated uncertainty.
\item \emph{Structural identifications}, ranked from
  strongest to weakest:
  \begin{itemize}
  \item[\textbf{(a)}] CMS $\to$ KZ $\to$ CS
    (proved operator isomorphism, EFK 1998);
  \item[\textbf{(b)}] dreibein $=$ linearised Havelock mode
    (proved at linearised level; nonlinear extension
    guaranteed by target-space symmetry of the WZW model);
  \item[\textbf{(c)}] Regge vertex $=$ point vortex
    (kinematic identity: deficit angle $=$ circulation;
    the Onsager selection on this system is the
    essential physical hypothesis,
    Lemma~\ref{III-lem:onsager-regge}).
  \end{itemize}
  Links~(a) and~(b) are mathematical isomorphisms.
  Link~(c)---the claim that ``the vortex system is
  spacetime''---is the theory's core physical
  hypothesis, experimentally supported (planetary polygons,
  BEC vortex experiments) and rests on the proved
  kinematic identity plus the Onsager variational
  principle.
\item \emph{Computed quantities} ($V_{cb}$, $H_0$, $\eta_B$,
  $\sin^2\theta_W$): the numerical values are exact consequences
  of the theory, but inherit the sensitivity of the structural
  identifications.
\end{enumerate}
The honest accounting is therefore: the theorem links are
rigorous; the structural identifications are the \emph{essential}
hypotheses; and the computed quantities test whether these
hypotheses are self-consistent.
The propagated uncertainty varies by orders of magnitude
across predictions, depending on the derivation chain:
\begin{center}\small
\begin{tabular}{llll}
\toprule
Prediction & Dominant source & Propagated & Match \\
\midrule
$\delta_{\mathrm{CKM}}$
  & radiative $O(c_{\mathrm{dn}}^2/c)$
  & $2.8\%$ ($2.0^\circ$) & $0.4\sigma$ \\
$\ln(M_P/v)$
  & Dunham $O(1/c_7^2)$
  & $0.04\%$ & $0.04\%$ \\
$H_0$
  & Dunham $O(1/c_{11}^2)$
  & $0.006\%$ & $< 0.1\%$ \\
$m_s/m_b$
  & Euler class $O(1/c_7^2)$
  & $0.05\%$ & $4\%$ \\
$\Omega_\Lambda/\Omega_{\mathrm{DM}}$
  & two-loop $O(1/c_{11})$
  & $0.8\%$ & $0.5\%$ \\
$|V_{cb}|$
  & $\delta\rho^*/\rho^* \times$ elast.\ $({-}2.5)$
  & $5\%$ & $6\%$ \\
$\sin^2\theta_W$
  & SM RG running to $M_Z$
  & ${\sim}4\%$ & $3.9\%$ \\
$\eta_B$
  & lattice QCD ($\kappa$, $\Gamma_y$)
  & factor~$3$ & factor~$3$ \\
\bottomrule
\end{tabular}
\end{center}
\noindent{\footnotesize The $\delta_{\mathrm{CKM}}$
uncertainty is the one-loop radiative correction to the
BF threshold $c_{\mathrm{dn}} = 3/2$: the Virasoro block
shift $\delta c = c_{\mathrm{dn}}^2/c = 0.044$ propagates
through the spectral-parameter integral as
$\delta\delta = (\pi/4)\tanh(\pi) \times \delta c = 2.0^\circ$,
which is $0.65\sigma$ of the experimental error ($\pm 3^\circ$).
The Harish-Chandra normalisation
(Paper~IV, Theorem~\ref*{IV-thm:ckm-phase}, Step~4)
is the standard spectral normalisation on~$\mathbf{H}^2$;
it is a structural identification, not a free parameter.}

\medskip
The first four predictions are in the \emph{algebraic tier}:
their uncertainties ($\le 0.8\%$) come from
$O(1/c^2)$ Dunham and Euler class corrections,
which decrease as~$N^{-4}$.
The next two ($|V_{cb}|$, $\sin^2\theta_W$) are in the
\emph{geometric tier}: their uncertainties ($4$--$5\%$)
reflect the sensitivity of radial wave functions and
RG running.
$\eta_B$ is limited by external (lattice QCD)
parameters, not by the theory.

\subsection{Structural limitations}
The geometric-tier predictions depend on fermion
zero-mode profiles on the Seifert manifold.
Including the $O(1/c)$ Euler class and BO corrections
(eq.~\eqref*{IV-eq:Dc-corrected}) brings $m_s/m_b$
from $0.045$ to $0.023$ ($4\%$ from observed),
promoting it to the algebraic tier.

Two predictions are beyond current experimental reach.
(i)~The frozen-mode dark matter has cross-section
$\sigma \sim 10^{-103}\;\mathrm{cm}^2$---undetectable by
any foreseeable direct-detection experiment.
The suppression is structural: $\sigma \propto e^{-c}$
with $c \approx 52$ (the topological coupling).
The prediction is testable indirectly through the
clustering spectrum: the frozen modes have a specific
mass spectrum (eq.~\eqref{V-eq:dm-masses}) that
imprints on small-scale structure.
(ii)~The extension to $d \ge 4$ spatial dimensions
uses polygons with $N \ge 10$ (unstable on the flat plane)
as \emph{mathematical probes} of the Green's function
angular structure.
The Havelock identity holds for all~$N$ regardless of
stability (Paper~I, Lemma~\ref*{I-lem:havelock});
the Onsager selection (H2b) applies only to the physical
$N = 7$ sector.
The $d \ge 4$ result is therefore a necessary condition
for Einstein manifolds (mode-independence of~$C_1$
requires spatial isotropy), not a full derivation---the
time-direction constraint requires the boundary Virasoro
mechanism (Paper~IV, \S\ref*{IV-sec:action}).

\subsection{Two precision regimes}
The predictions fall into two tiers by mechanism,
not by post-hoc selection:
\begin{center}
\small
\begin{tabular}{lll}
\toprule
Tier & Mechanism & Accuracy \\
\midrule
\emph{Algebraic} &
  Eigenvalue sums, phases, winding numbers &
  $< 5\%$ \\
& ($\delta_{\mathrm{CKM}}$, $|V_{cb}|$, $H_0$,
  $\Omega_\Lambda$, $\Omega_{\mathrm{DM}}$, $\Omega_b$,
  $\Delta m^2$ ratio, $m_H$, $m_s/m_b$) & \\[3pt]
\emph{Geometric} &
  Radial profiles on $\mathbf{H}^2$ &
  factor~${\sim}2$ \\
& ($m_c/m_t$, $m_u$, $m_d/m_s$,
  $|V_{us}|$, $\sin^2\theta_W$) & \\
\bottomrule
\end{tabular}
\end{center}
\noindent
The algebraic tier uses the exact Havelock decomposition
($\lambda_m = C_1 - m(N{-}m)/2$, a theorem);
the geometric tier uses the Dirac zero-mode profile
$\psi_c(\rho) = (\sinh\rho)^{c-1/2}$ on the Seifert manifold
(zero adjustable parameters per generation, versus
${\sim}\!21$ in standard RS models).
The pattern is structural, not accidental: it reflects
the separation between topological data (exact)
and radial geometry (approximate).

\subsection{Comparison to zero-parameter frameworks}
\label{sec:comparison}
The theory's accuracy should be compared not to the
Standard Model (which fits $19$+ parameters by construction)
but to other frameworks that predict fermion masses
from structure alone:
\begin{itemize}
\item Tree-level SU(5) (Georgi--Glashow 1974):
  $\sin^2\theta_W = 3/8$ ($62\%$ off);
  $m_b/m_\tau = 1$ (factor $2.4$ off).
  Requires MSSM to reach $7\%$ accuracy.
\item Chiral perturbation theory (NLO, $12$ parameters):
  $K \to \pi\pi$ amplitudes at $20$--$30\%$.
\item Standard RS (anarchic Yukawas, $21$+ parameters):
  $< 1\%$ accuracy, but every mass is fitted, not predicted.
\item String landscape (Denef--Douglas 2004):
  $\sim\!10^{500}$ vacua; SM parameters are environmental,
  not computed. No single vacuum is identified.
\item Asymptotic safety (Eichhorn--Held 2018):
  $\sin^2\theta_W$ from UV fixed-point constraints;
  fermion mass predictions are at the one-loop level only.
\end{itemize}
Among frameworks with zero continuous free parameters,
no other simultaneously predicts
fermion mass ratios, CKM elements,
and cosmological parameters from a single discrete input.
The algebraic-tier predictions ($< 5\%$) include $m_s/m_b$,
which reaches $4\%$ accuracy once the $O(1/c)$ Euler class
correction is included.

\subsection{Falsifiable predictions}
\label{sec:falsifiable}

\noindent
The central-value matches on~$V_{cb}$
($\pm 12\%$ theoretical sensitivity, elasticity $-2.5$) and~$H_0$
($\pm O(1/c^2)$ from BO truncation) should be read as
\emph{consistency} of the theory at the level of its
structural identifications, not as independently
predicted values.  The genuine predictions---those that
could \emph{falsify} the theory---are of three kinds.
\emph{(i)~Binary tests}
(imminent: JUNO for the neutrino hierarchy,
DESI/Euclid for $w = -1$,
LEGEND/nEXO for the absence of neutrinoless double beta decay):
\begin{itemize}
\item Normal neutrino hierarchy (if inverted $\to$ falsified);
\item $w = -1$ exactly, at all redshifts
  (the tree-level free energy $F_{\mathrm{DE}}$ is
  $\rho$-independent: a cosmological constant with no
  dynamical component).
  The DESI DR2 BAO data (2025) show a
  $2.8$--$4.2\sigma$ preference for time-varying $w$
  in the $w_0 w_a$CDM parametrisation;
  the best-fit values ($w_0 \approx 0$, $w_a \approx -3.7$)
  require phantom crossing, which no standard scalar field
  produces.
  If $w \ne -1$ is confirmed at $5\sigma$ with independent
  datasets, the theory is falsified;
\item No $0\nu\beta\beta$ (Dirac neutrinos from the Seifert structure:
  the neutrino masses arise from the Dirac seesaw with a KK mass
  on~$S^1$, not a Majorana mass, so lepton number is exactly conserved);
\item Proton exactly stable ($\mathbb{Z}_{56}$ conservation;
  Hyper-K will reach $\tau > 10^{35}$\,yr);
\item $\theta_{\mathrm{QCD}} = 0$ exactly
  (CS level quantisation for $\mathrm{SU}(3)$,
  preserved under fermion shifts;
  no axion, no axion dark matter;
  nEDM $= 0$ to all orders; ADMX null result).
\item EWPT gravitational wave signal at
  $f \approx 4 \times 10^{-5}$\,Hz,
  $\Omega_{\mathrm{GW}} h^2 \approx 10^{-9}$
  (eq.~\eqref{V-eq:gw-ewpt};
  the SM gives $\Omega = 0$; LISA $\sim\!2037$).
\item $\Omega_k = 0$ exactly
  (the KK reduction absorbs all spatial curvature into
  $\Lambda_3 = (N^2{-}16)/16$;
  CMB-S4 and Euclid will reach $|\Omega_k| < 10^{-4}$,
  where the theory predicts exactly zero,
  distinguishable from inflationary
  $\Omega_k \sim e^{-2N_e}$;
  \S\ref{V-sec:geometry-initial}).
\end{itemize}
\emph{(i\hspace{0.5pt}b)~Null predictions}
(the absence of each is a falsifiable consequence):
no cosmic strings
(the gauge group $\mathrm{SU}(3) \times \mathrm{SU}(2) \times \mathrm{U}(1)$
is simply connected in the covering group; the discrete
$\mathbb{Z}_N$ symmetry breaking at the polygon transition
produces domain walls, not strings, and these are erased by
inflation),
no magnetic monopoles
(no GUT-scale symmetry breaking; the gauge group
is built from $k = 1$ CS sectors, not broken from a
larger group),
no isocurvature perturbations (single breathing mode
$\rho$ provides a single clock),
$N_{\mathrm{eff}} = 3.00$ (no additional light relics),
no cosmic birefringence,
no dark matter annihilation signals
(gravitational-strength coupling,
$\sigma_{\mathrm{SI}} \sim 10^{-103}\;\mathrm{cm^2}$),
NANOGrav stochastic background is astrophysical
(supermassive black hole binaries, not cosmological),
and $r \sim 10^{-6}$ (undetectable primordial gravitational
waves from inflation; \S\ref{V-sec:inflation}).

\emph{(i\hspace{0.5pt}c)~Planetary polygon predictions}
(Paper~II, \S\ref*{II-sec:oblateness}):
\begin{itemize}
\item Saturn hexagon stability requires oblateness
  ($f = 0.098$): on a mean sphere, $\lambda_3 = -0.12$
  (unstable); on the oblate spheroid, $\lambda_3 = +0.01$
  (marginally stable).
  The critical latitude is $72^\circ$N; the hexagon
  at $78^\circ$N is $6^\circ$ above the boundary.
  This is confirmed by Cassini data (40-year persistence,
  deep barotropic column).
\item Neptune: $N \in \{3, 4, 5\}$ (packing-dominated;
  $\beta R^3/\kappa \approx 0.1 \ll 1$).
  Testable by a Neptune polar orbiter at $\sim 100$\,km
  resolution.
  $N \ge 6$ would falsify the packing bound.
\item Uranus: $N = 0$ (no polygon; $98^\circ$ axial tilt
  disrupts Onsager condensation).
  A polar polygon would challenge the tilt-disruption
  hypothesis.
  Current data ($N = 1$, single cyclone;
  \citealt{Akins2023}) is consistent.
\end{itemize}

\emph{(i\hspace{0.5pt}d)~Confinement and $\alpha_s$.}
The polygon phase confines
(Proposition~\ref{IV-prop:confinement}):
$\gcd(7,3) = 1$ forces the $\mathbb{Z}_3$ center symmetry
of $\mathrm{SU}(3)$ to be unbroken, giving
$\langle P\rangle = 0$.
The deconfinement transition is the polygon--BTZ
Hawking--Page transition.
The strong coupling $\alpha_s(M_{\mathrm{poly}})
= 1/(4\pi\sqrt{2}) = 0.0563$
(eq.~\eqref{IV-eq:alpha-s}) is a zero-parameter
prediction; one-loop running gives
$\alpha_s(M_Z) = 0.115$,
two-loop gives $0.122$, bracketing the
observed $0.1180 \pm 0.0009$.
The QCD string tension
$\sigma/\Lambda^2 = 5.97$
(Derivation~\ref{IV-der:string-tension})
matches the lattice value $6.25 \pm 0.5$
(\citealt{Bali2001}) to~$5\%$: the four ingredients
(Havelock eigenvalue, Polyakov monopole gas on~$\mathbf{H}^2$
with $\mathrm{rank}(\mathrm{SU}(3)) = 2$ species,
Verlinde dimension, and CS partition function) are each
computed from~$N$ and~$k$ with no free parameters.

\emph{(ii)~Correlated observables.}
The instanton fugacity $K = 0.548$
(eq.~\eqref{IV-eq:instanton-K},
determined by $b(N)$ with no free parameters;
the $4\%$ match on $m_s/m_b$
confirms that $K$ is computed consistently).
$K$~links two independent observables:
$m_s/m_b \propto 1/K^2$ (eq.~\eqref{V-eq:mass-hierarchy},
elasticity $\partial\ln(m_s/m_b)/\partial k_{\mathrm{frac}} = 4\pi$)
and $\Omega_b$ via
$K/(1{+}K^2)$ (eq.~\eqref{V-eq:mass-gap-resum},
elasticity $-0.3$).
The CKM phase $\delta_{\mathrm{CKM}} = 70.2^\circ$
is a \emph{separate} prediction, independent of~$K$
(it depends on the spectral parameter $s = 1$ at the
BF crossing, not on the CS level).
In the Standard Model, $m_s/m_b$ and $\Omega_b$ are unrelated;
here a single computed number~$K$ determines both.
Any precision measurement that breaks the $K$-correlation
falsifies the theory.
\emph{(iii)~Structural derivations
that the SM treats as inputs}: the gauge group,
$3$~generations, $m_u < m_d$, $\theta_{\mathrm{QCD}} = 0$.

\noindent
The perturbative CKM gives $J = 3.5 \times 10^{-5}$
(observed $3.0 \times 10^{-5}$, $17\%$~match)---a genuine
prediction with all entries computed from~$N = 7$.

\paragraph{Cosmological structure.}
The theory determines the large-scale structure of the
universe without free cosmological parameters beyond~$M_P$.
Spatial sections are flat
($\Omega_k = 0$, \S\ref{V-sec:geometry-initial}),
the equation of state is $w = -1$ at all redshifts
(\S\ref{V-sec:cc-instanton}),
and the initial state is uniquely determined by the
Born--Oppenheimer structure of the polygon breathing mode
(Remark~\ref{III-rmk:wdw-uniqueness}).
The classical big-bang singularity is replaced by quantum
tunnelling at the palindromic threshold~$\rho^*$, where
the 2D vortex system on~$\mathbf{H}^2$ gives rise to
$3{+}1$-dimensional Lorentzian spacetime through the
Wick rotation $\rho \to \rho + i\pi/2$.
The Bekenstein--Hawking entropy at the threshold
is determined by the same central charge
$c = 12\,b(N)$ that governs the WDW kinetic term
(Remark~\ref{III-rmk:entropy-bridge}):
$S_{\mathrm{BH}} = (\pi c/3)\cosh\rho^*$, counting the
non-perturbative Virasoro descendants of the angular modes
coarse-grained by the Onsager mechanism.

\paragraph{Why $3{+}1$ dimensions.}
The polynomial Havelock eigenvalue $\lambda_m = (N{-}1) - m(N{-}m)/2$
is unique to two spatial dimensions
(Remark~\ref{I-rmk:polynomial-uniqueness}):
the Ramanujan sum identity
$\sum[1{-}\cos(2\pi pm/N)]/\sin^2(\pi p/N) = 2m(N{-}m)$
has no analogue for $\sin^\alpha$ with $\alpha \ne 2$,
and the angular Hessian kernel in $d$~dimensions
involves $\sin^{-d}$.
The polynomial structure enables CMS integrability
(from the holomorphicity of $\log z$), exact BO
separation, and the entire Havelock framework.
The third spatial dimension arises from the $S^1$ fibre
of the Seifert manifold ($\mathbb{Z}_N \subset U(1)$
is the polygon rotational symmetry).
Time emerges from the WKB approximation on the single
breathing mode~$\rho$---the unique slow degree of
freedom after the CMS angular modes are frozen.
The count $2\,(\text{log kernel}) + 1\,(S^1\text{ fibre})
+ 1\,(\text{WKB clock}) = 3{+}1$ is forced by
the Havelock eigenvalue structure.
The WKB time coordinate is valid in the semiclassical
regime ($\rho > \rho^*$, $c \gg 1$); at the turning
point $\rho = \rho^*$ the Airy regime replaces it,
and for $\rho < \rho^*$ there is no classical time
(\S\ref{V-sec:geometry-initial}).

\paragraph{Structure formation and the high-redshift universe.}
The radion inflection-point potential
(\S\ref{V-sec:inflation}) gives a featureless red tilt
with $n_s(k) \approx 0.964$--$0.970$ across galactic scales
($k = 1$--$100\;\mathrm{Mpc}^{-1}$) and running
$dn_s/d\ln k \approx -10^{-3}$.
The four dark matter species (mass ratios
$1{:}\sqrt{3}{:}\sqrt{6}{:}\sqrt{10}$, equal energy density
per species, $\sigma_{\mathrm{SI}} \sim 10^{-103}\;\mathrm{cm}^2$)
produce standard NFW haloes: cold, collisionless, and
gravitationally indistinguishable from single-component CDM
on all astrophysical scales.
The slightly higher $n_s = 0.970$ (relative to the
Planck central value $0.965$) increases the halo abundance
at $z > 10$ by $\sim\!10$--$20\%$; the remainder of any
high-redshift galaxy excess is accounted for by the
well-established increase in star formation efficiency
at low metallicity: reduced feedback
(metal-poor stellar winds), gas-rich haloes
(no prior depletion), top-heavy initial mass functions
(larger Jeans mass $\to$ more UV per baryon),
and bursty star formation histories
that bias UV-selected samples upward.
JWST observations at $z > 10$
\citep{Labbe2023,Castellano2022,Naidu2022}
are quantitatively consistent with the theory's
power spectrum provided $\varepsilon_{\mathrm{SF}} \approx
0.15$--$0.25$ (compared to $\varepsilon \approx 0.05$ locally),
a range supported by high-resolution simulations
of metal-free star formation \citep{Bromm2013}.
The theory predicts that future measurements of the
matter power spectrum at $k > 1\;\mathrm{Mpc}^{-1}$
(Euclid, Roman) will find $n_s \approx 0.96$--$0.97$
with no blue tilt or scale-dependent enhancement.

\section{Conclusion}
\label{sec:conclusion}

The audit reveals a sharp dichotomy.
The algebraic predictions---CKM phase, Hubble parameter,
neutrino mass ratio---match data to better than~$5\%$
through exact eigenvalue arithmetic.
The geometric-tier predictions ($m_c/m_t$, $m_u$,
$|V_{us}|$, $\sin^2\theta_W$) remain at factor-${\sim}2$
accuracy, reflecting the sensitivity of light-quark masses
to the Dirac zero-mode profiles on~$\mathbf{H}^2$.

Seven binary tests are experimentally
accessible within two decades.
Of these, the EWPT gravitational wave signal
($f \approx 4 \times 10^{-5}$\,Hz,
$\Omega_{\mathrm{GW}} h^2 \approx 10^{-9}$),
the $K$-correlation between $m_s/m_b$ and $\Omega_b$,
and $\Omega_k = 0$ exactly are \emph{diagnostic}:
they distinguish this theory from other BSM approaches.
The remaining four (normal hierarchy, $w = -1$, proton
stability, $\theta_{\mathrm{QCD}} = 0$) are shared
with the SM or $\Lambda$CDM.
A single failure of a binary test (normal hierarchy,
$w = -1$, proton stability, $\theta_{\mathrm{QCD}} = 0$)
falsifies the theory;
failure of a quantitative prediction ($\sin^2\theta_W$,
$|V_{cb}|$, $\eta_B$) constrains the derivation chain
at the failed step.
A companion supplement develops an arithmetic interpretation:
the Havelock eigenvalues satisfy the Ramanujan--Petersson
bound for all stable polygons, linking vortex stability
to Langlands temperedness.

\bibliographystyle{plainnat}
\bibliography{../shared/refs}

\end{document}
